\section{数据 Scaling Laws：理论与经验}

\subsection{Kaplan 论文的经验观察}

Kaplan et al. 的 Scaling Laws 论文是该领域最全面的参考文献之一。其核心经验发现是：\textbf{在 log-log 图上，测试损失与计算量、数据量、参数量之间呈线性关系}——即存在幂律关系。

\begin{figure}[H]
\centering
\includegraphics[width=0.95\textwidth]{slides-images/slide-007.jpg}
\caption{Kaplan et al.：计算、数据量、参数量与测试损失的幂律关系\protect\footnotemark}
\end{figure}
\footnotetext{Slides 第8页。}

\begin{importantbox}{幂律关系的数学表达}
当在 log-log 图上观察到线性关系时，意味着：
$$\log L = -\alpha \log X + c$$
即 $L = c' \cdot X^{-\alpha}$，其中 $L$ 是测试损失，$X$ 是资源量（数据/参数/计算），$\alpha > 0$ 是缩放指数。
\end{importantbox}

一个重要的细微之处是：当我们单独缩放某个变量（如数据量）时，\textbf{其他变量必须足够大}，以避免进入渐近饱和区域。例如，研究数据缩放时，模型尺寸必须大到不会被数据``填满''。

\subsection{为什么幂律衰减是自然的？——均值估计}

Tatsu 从最简单的统计问题出发，解释为什么幂律衰减是``自然的''。

\textbf{例子 1：估计高斯均值}

给定 $n$ 个独立同分布样本 $X_1, \ldots, X_n \sim \mathcal{N}(\mu, \sigma^2)$，用样本均值 $\hat{\mu} = \frac{1}{n}\sum X_i$ 估计 $\mu$。

估计误差为：
$$\mathbb{E}[(\hat{\mu} - \mu)^2] = \frac{\sigma^2}{n}$$

取对数后：
$$\log(\text{Error}) = 2\log\sigma - \log n$$

\begin{figure}[H]
\centering
\includegraphics[width=0.95\textwidth]{slides-images/slide-009.jpg}
\caption{均值估计的 Scaling Law：误差以 $n^{-1}$ 的速率衰减\protect\footnotemark}
\end{figure}
\footnotetext{Slides 第10页。}

这就是最简单的 Scaling Law：在 log-log 图上斜率为 $-1$。

\subsection{为什么实际观测到的指数很小？——非参数回归与维度诅咒}

然而，实际观察到的缩放指数远小于理论预期。例如：
\begin{itemize}
    \item 机器翻译（Hestness）：$\alpha = 0.13$
    \item 语音识别（Hestness）：$\alpha = 0.3$
    \item 语言建模（Kaplan）：$\alpha \approx 0.095$
\end{itemize}

这些指数远远慢于简单估计问题的 $\alpha = 1$ 或 $\alpha = 0.5$。为什么？

\begin{figure}[H]
\centering
\includegraphics[width=0.95\textwidth]{slides-images/slide-010.jpg}
\caption{实际观测的缩放指数远小于经典理论预期\protect\footnotemark}
\end{figure}
\footnotetext{Slides 第11页。}

\textbf{例子 2：非参数回归与维度诅咒}

考虑在二维空间中估计任意回归函数 $f$。最简单的方法是将空间切分成小方块，在每个方块内取均值。

\begin{figure}[H]
\centering
\includegraphics[width=0.95\textwidth]{slides-images/slide-011.jpg}
\caption{非参数回归的直觉：将空间切分成小方块，在每个方块内进行均值估计\protect\footnotemark}
\end{figure}
\footnotetext{Slides 第12页。}

如果有 $n$ 个样本在 $d$ 维空间中：
\begin{itemize}
    \item 将空间切成 $n^{1/2}$ 个方块（二维）$\rightarrow$ 每个方块分到 $n^{1/2}$ 个样本 $\rightarrow$ 误差 $\sim n^{-1/2}$
    \item 推广到 $d$ 维：误差 $\sim n^{-1/d}$
\end{itemize}

$$\text{Error} \sim n^{-1/d}$$

在 log-log 图上，斜率为 $-1/d$。

\begin{importantbox}{缩放指数反映内在维度}
缩放指数 $\alpha = 1/d$ 直接与数据的\textbf{内在维度}（intrinsic dimensionality）相关。语言建模的缩放指数 $\alpha \approx 0.095$ 意味着学习任务的有效维度大约为 $d \approx 10.5$。换言之，缩放越慢，说明任务越``复杂''（有效维度越高）。
\end{importantbox}

\begin{figure}[H]
\centering
\includegraphics[width=0.95\textwidth]{slides-images/slide-012.jpg}
\caption{内在维度与缩放指数的理论和经验对比\protect\footnotemark}
\end{figure}
\footnotetext{Slides 第13页。}

\begin{warningbox}{内在维度估计的困难}
虽然理论上缩放指数与内在维度密切相关，但\textbf{直接估计数据的内在维度}是一个极其困难的问题——其难度与建模数据本身相当。因此，这种对应关系更多是概念性的，不应过度解读。
\end{warningbox}

\subsection{数据 Scaling Laws 的实用价值}

数据 Scaling Laws 不仅仅是理论上的优美结论，它还有重要的工程应用。

\begin{figure}[H]
\centering
\includegraphics[width=0.95\textwidth]{slides-images/slide-013.jpg}
\caption{数据组成只影响偏移量不影响斜率——可在小规模上做数据选择\protect\footnotemark}
\end{figure}
\footnotetext{Slides 第14页。}

\textbf{应用 1：数据混合优化}

Kaplan 发现数据组成（如 Wikipedia vs. Common Crawl 的比例）\textbf{只影响 Scaling Law 的偏移量，不影响斜率}。这意味着你可以在小模型上做数据选择实验，结论直接迁移到大模型。

\textbf{应用 2：多轮训练的衰减}

\begin{figure}[H]
\centering
\includegraphics[width=0.95\textwidth]{slides-images/slide-014.jpg}
\caption{多 Epoch 训练的有效样本量递减：约 4 个 epoch 后收益急剧下降\protect\footnotemark}
\end{figure}
\footnotetext{Slides 第15页。}

``我们快要用完互联网数据了''的担忧引出了另一个重要问题：重复训练同一批数据的效果如何？研究表明存在一个\textbf{有效样本量}的概念：约 4 个 epoch 之后，重复数据的收益急剧递减。

\textbf{应用 3：重复高质量数据 vs. 引入低质量新数据}

\begin{figure}[H]
\centering
\includegraphics[width=0.95\textwidth]{slides-images/slide-015.jpg}
\caption{数据重复 vs. 引入新数据的权衡\protect\footnotemark}
\end{figure}
\footnotetext{Slides 第16页。}

当你需要数万亿 token 进行训练时，面临选择：是重复 Wikipedia 等高质量数据 10 次，还是引入更多低质量但全新的数据？CMU 的研究者们利用 Scaling Laws 框架来分析这种权衡——假设每种数据源都有独立的幂律关系，通过在小规模上拟合每种混合比例的 Scaling Law，就能预测大规模下的最优混合策略。

\subsection{本章小结}

数据 Scaling Laws 是所有 Scaling Laws 中最直观、理论基础最扎实的。幂律衰减可以从简单的统计问题中自然推导出来，而实际观测到的缓慢指数反映了学习任务的内在复杂度。这些规律不仅在经验上极其稳健（跨领域、跨模型族成立），还有重要的工程应用价值。

%% ========== Part 4: 模型 Scaling ==========

